%% LyX 2.4.4 created this file.  For more info, see https://www.lyx.org/.
%% Do not edit unless you really know what you are doing.
\documentclass[english]{article}
\usepackage[T1]{fontenc}
\usepackage[latin9]{inputenc}
\usepackage{units}
\usepackage{enumitem}
\usepackage{amsmath}
\usepackage{amsthm}
\usepackage{geometry}
\geometry{verbose,tmargin=1in,bmargin=1in,lmargin=1in,rmargin=1in}

\makeatletter
%%%%%%%%%%%%%%%%%%%%%%%%%%%%%% Textclass specific LaTeX commands.
\numberwithin{equation}{section}
\numberwithin{figure}{section}
\newlength{\lyxlabelwidth}      % auxiliary length 

%%%%%%%%%%%%%%%%%%%%%%%%%%%%%% User specified LaTeX commands.
\usepackage{fancyhdr}
\usepackage{lastpage}
\pagestyle{fancy}
\usepackage{enumitem}
\usepackage{ifthen}
%sets math fonts to be more readable
\everymath=\expandafter{\the\everymath\displaystyle}
%\DeclareMathSizes{10}{10}{10}{10}
\usepackage{multicol}

\usepackage{tikz}
\usetikzlibrary{plotmarks}
\usepackage{pgfplots}
\usetikzlibrary{shapes.geometric, intersections}
\pgfplotsset{compat=newest}

\usepackage{calc}
\usepackage[nomessages]{fp}

\usepackage{advdate}    % Advancing/saving dates
\usepackage{datetime}   % Dates formatting
\usepackage{datenumber} % Counters for dates
% Added by lyx2lyx
\setlength{\parskip}{\smallskipamount}
\setlength{\parindent}{0pt}

\makeatother

\usepackage{babel}
\begin{document}
%%Header Info
\renewcommand{\author}{Dr.\,James Ashe}
\newcommand{\classNum}{137}
\newcommand{\classSec}{B}
\newcommand{\nameType}{Name:}
\newcommand{\examType}{Exam 3}
\renewcommand{\date}{\formatdate{15}{4}{2013}}%%First page's header%%%%%%%%%%%%%%%%%%%%%%
\fancypagestyle{firststyle} %%header settings for first pagr
{
	\fancyhead[L]{MTH \classNum{}(\classSec{}) \author{}\\
	\textbf{\examType{}} \date{}}
	\fancyhead[C]{\textbf{\nameType}}
	\setlength{\headsep}{14pt}
}
%%%%%%%%%%%%%%%%%%%%%%%%%%%%%%%%%%%%%%%%%%%%Next page's header%%%%%%%%%%%%%%%%%%%%%%%
\setlist{nolistsep}
\lhead{\classNum{}(\classSec{})} 
\chead{\textbf{\nameType}} 
\cfoot{\thepage{}~of~\pageref{LastPage}} 
\setlength{\headsep}{5pt} 
%%%%%%%%%%%%%%%%%%%%%%%%%%%%%%%%%%%%%%%%%%%%%Various settings; see the preamble for other settings%%%%%
%\setenumerate{itemsep=.125cm,leftmargin=12pt,itemindent=0pt,labelwidth=0pt}
\setlist{itemsep=.05cm}
\renewcommand{\labelenumi}{\textbf{\arabic{enumi}.}}
\renewcommand{\labelenumii}{\textbf{\alph{enumii}.}}
\thispagestyle{firststyle}
\newcommand{\dspace}{\vspace*{.75cm}}
%%%%%%%%%%%%%%%%%%%%%%%%%%%%%%%%%%%%%%%%%%
\newcommand{\xmax}{0}
\newcommand{\xmin}{-\xmax}
\newcommand{\xstep}{0}
\newcommand{\ymax}{0}
\newcommand{\ymin}{-\ymax}
\newcommand{\ystep}{0} 

\textbf{You must show your work to receive credit. }Unless stated
otherwise, all problems are worth 5 points. 
\begin{enumerate}[resume]
\item {[}2pts each{]}\label{enu:quad-graph} Use the graph of the quadratic
function, to answer the following:

\begin{multicols}{2}
\begin{enumerate}
\item What is the vertex?\dspace
\item What is the axis of symmetry?\dspace
\item What is the range?\dspace
\item What is the domain?\dspace
\item Using the graph, what interval(s) is the function positive? (approximate)\dspace\columnbreak
\item []\renewcommand{\xmax}{10}
\renewcommand{\xmin}{-10}
\renewcommand{\xstep}{0} %must be xmin+n
\renewcommand{\ymax}{10}
\renewcommand{\ymin}{-10}
\renewcommand{\ystep}{0} %must be ymin+n
%%%%%%%
\begin{tikzpicture} 
\begin{axis}[height=5.5cm, 
	width=5.5cm, 
	axis lines=middle,
	minor x tick num={4},
	minor y tick num={4},
	grid=both,
	%xtick={0,50,...,250},
	%ytick={0,10,20,...,40},
	%axis line style={-latex}, 
	xmin=\xmin,xmax=\xmax, 
	ymin=\ymin,ymax=\ymax,
	%restrict y to domain=-1:10,
	%no markers, 
	xlabel={$x$},ylabel={$y$}, 
	%every axis x label/.append style={anchor=north}, 
	%every axis y label/.append style={anchor=east}
	]
	
	%%\addplot[color=red,solid,mark=*] coordinates {(-3,-4)};
	\addplot[domain=\xmin:\xmax,thick,samples=100,draw=red,<->]{-(x-2)^2+5} [yshift=8pt] node[pos=0.7,right] {$g(x)$};
	
\end{axis}
\end{tikzpicture}
\end{enumerate}
\end{multicols}
\item Using the function $f(x)=x^{2}-4x+4$, complete the following:

\begin{multicols}{2}
\begin{enumerate}
\item {[}2pts{]} Determine $f$'s axis of symmetry.\dspace
\item {[}2pts{]} If any exist, find the $x$-int(s).\dspace
\item {[}2pts{]} Find the $y$-int.\dspace
\item {[}4pts{]} Sketch the graph\columnbreak
\item []\renewcommand{\xmax}{10}
\renewcommand{\xmin}{-10}
\renewcommand{\xstep}{0} %must be xmin+n
\renewcommand{\ymax}{10}
\renewcommand{\ymin}{-10}
\renewcommand{\ystep}{0} %must be ymin+n
%%%%%%%
\begin{tikzpicture} 
\begin{axis}[height=5.5cm, 
	width=5.5cm, 
	axis lines=middle,
	minor x tick num={4},
	minor y tick num={4},
	grid=both,
	%xtick={0,50,...,250},
	%ytick={0,10,20,...,40},
	%axis line style={-latex}, 
	xmin=\xmin,xmax=\xmax, 
	ymin=\ymin,ymax=\ymax,
	%restrict y to domain=-1:10,
	no markers, 
	xlabel={$x$},ylabel={$y$}, 
	%every axis x label/.append style={anchor=north}, 
	%every axis y label/.append style={anchor=east}
	]
	
	%%\addplot[color=red,solid,mark=*] coordinates {(-3,-4)};
	%%\addplot[domain=\xmin:\xmax,thick,samples=100,draw=red,<->]{-(x-2)^2+5} node[right] {$$};
	
\end{axis}
\end{tikzpicture}
\end{enumerate}
\end{multicols}
\item Consider $P(x)=x^{2}-5x+2$ and $(x-2)$.

\begin{enumerate}
\item {[}4pts{]} Use division (synthetic or traditional) to find the quotient
and remainder when $P(x)$ is divided by $x-2$. \vfill\dspace
\item {[}4pts{]} Find $\frac{P(x)}{x-2}$\vfill
\item {[}2pts{]} Is $x=2$ a root/zero of $P(x)$?\dspace\newpage
\end{enumerate}
\item Consider the polynomial $f(x)$ shown in the graph.

\begin{multicols}{2}
\begin{enumerate}
\item What are the zeros of $f(x)$?\dspace\dspace
\item Find the simplest polynomial with real coefficients that has the roots
from above (you do not have to expand).\dspace\columnbreak
\item []\renewcommand{\xmax}{7}
\renewcommand{\xmin}{-5}
\renewcommand{\xstep}{0} %must be xmin+n
\renewcommand{\ymax}{15}
\renewcommand{\ymin}{-25}
\renewcommand{\ystep}{0} %must be ymin+n
%%%%%%%
\begin{tikzpicture} 
\begin{axis}[height=5.5cm, 
	width=5.5cm, 
	axis lines=middle,
	minor x tick num={4},
	minor y tick num={4},
	grid=both,
	%xtick={0,50,...,250},
	%ytick={0,10,20,...,40},
	%axis line style={-latex}, 
	xmin=\xmin,xmax=\xmax, 
	ymin=\ymin,ymax=\ymax,
	%restrict y to domain=-1:10,
	no markers, 
	xlabel={$x$},ylabel={$y$}, 
	%every axis x label/.append style={anchor=north}, 
	%every axis y label/.append style={anchor=east}
	]
	
	%%\addplot[color=red,solid,mark=*] coordinates {(-3,-4)};
	\addplot[domain=\xmin:\xmax,thick,samples=100,draw=red,<->]{-(x+2)*(x-2)*(x-5)} node[pos=.65,right] {$f(x)$};
	
\end{axis}
\end{tikzpicture}
\end{enumerate}
\end{multicols}
\item Use $f(x)=x^{3}-4x$ to complete the following.

\begin{multicols}{2}
\begin{enumerate}
\item {[}3pts{]} Find the $x$-int(s).\dspace
\item {[}3pts{]} Determine if the graph crosses or touches without crossing
the $x$-axis at each $x$-int.\dspace
\item {[}4pts{]} Sketch a graph of the function.\columnbreak
\item []\renewcommand{\xmax}{8}
\renewcommand{\xmin}{-5}
\renewcommand{\xstep}{0} %must be xmin+n
\renewcommand{\ymax}{10}
\renewcommand{\ymin}{-10}
\renewcommand{\ystep}{0} %must be ymin+n
%%%%%%%
\begin{tikzpicture} 
\begin{axis}[height=5.5cm, 
	width=5.5cm, 
	axis lines=middle,
	minor x tick num={4},
	minor y tick num={4},
	grid=both,
	%xtick={0,50,...,250},
	%ytick={0,10,20,...,40},
	%axis line style={-latex}, 
	xmin=\xmin,xmax=\xmax, 
	ymin=\ymin,ymax=\ymax,
	%restrict y to domain=-1:10,
	no markers, 
	xlabel={$x$},ylabel={$y$}, 
	%every axis x label/.append style={anchor=north}, 
	%every axis y label/.append style={anchor=east}
	]
	
	%%\addplot[color=red,solid,mark=*] coordinates {(-3,-4)};
	%%\addplot[domain=\xmin:\xmax,thick,samples=100,draw=red,<->]{-(x-2)^2+5} node[right] {$$};
	
\end{axis}
\end{tikzpicture}
\end{enumerate}
\end{multicols}
\item {[}5pts each{]} Find all real solutions.

\begin{enumerate}
\item $\sqrt{x+3}=x+1$\vfill\dspace
\item $x^{-\nicefrac{2}{5}}=4$\vfill
\end{enumerate}
\item {[}5pts each{]} Find the vertex of the graph of the following functions:

\begin{enumerate}
\item $f(x)=x^{2}-4x+10$.\dspace\dspace\dspace
\item $f(x)=-3\left(x+2\right)^{2}-5$.\dspace
\end{enumerate}
\end{enumerate}
\newpage\renewcommand{\dspace}{\vspace*{1.25cm}}
\begin{enumerate}[resume]
\item Use $f(x)=\frac{2x+6}{x-2}$

\begin{multicols}{2}
\begin{enumerate}
\item {[}4pts{]} Find the domain of $f(x)$.\dspace
\item {[}4pts{]} Find all vertical asymptotes.\dspace
\item {[}6pts{]} Find the horizontal asymptote.\dspace\dspace
\item {[}6pts{]} Sketch a graph of the function\columnbreak
\item []\renewcommand{\xmax}{10}
\renewcommand{\xmin}{-10}
\renewcommand{\xstep}{0} %must be xmin+n
\renewcommand{\ymax}{10}
\renewcommand{\ymin}{-10}
\renewcommand{\ystep}{0} %must be ymin+n
%%%%%%%
\begin{tikzpicture} 
\begin{axis}[height=5.5cm, 
	width=5.5cm, 
	axis lines=middle,
	minor x tick num={4},
	minor y tick num={4},
	grid=both,
	%xtick={0,50,...,250},
	%ytick={0,10,20,...,40},
	%axis line style={-latex}, 
	xmin=\xmin,xmax=\xmax, 
	ymin=\ymin,ymax=\ymax,
	%restrict y to domain=-1:10,
	no markers, 
	xlabel={$x$},ylabel={$y$}, 
	%every axis x label/.append style={anchor=north}, 
	%every axis y label/.append style={anchor=east}
	]
	
	%%\addplot[color=red,solid,mark=*] coordinates {(-3,-4)};
	%%\addplot[domain=\xmin:\xmax,thick,samples=100,draw=red,<->]{-(x-2)^2+5} node[right] {$$};
	
\end{axis}
\end{tikzpicture}
\end{enumerate}
\end{multicols}
\item {[}2pts each{]} Find the following:

\begin{enumerate}
\item $x^{2}-20$ as $x\rightarrow-\infty$\dspace
\item $-2x^{13}+x^{4}-50$ as $x\rightarrow\infty$\dspace
\item ${\displaystyle \lim_{x\rightarrow\infty}\frac{1}{x^{3}}}$\dspace
\item ${\displaystyle \lim_{x\rightarrow0^{+}}\frac{3}{x}}$\dspace
\item ${\displaystyle \lim_{x\rightarrow2^{-}}\frac{3}{x-2}}$\dspace
\item The domain of $\frac{x^{2}-3x-7}{x(x-2)(x+1)}$\dspace
\end{enumerate}
\end{enumerate}

\end{document}
